% !TeX root = ../教材题库_全部条目_转写工作稿.tex
% Source: XJ-B1-C06-HW-P256-Y0708
% 第6章 / 6.4 用样本估计总体 / 教材 p.256 / PDF 第263页
% OCR未校对：公式与文字均可能有误；未完成数学TeX转写；禁止直接发给学生
% 原题及相邻原解截图（核对用）：../题目原图/XJ-B1-C06-HW-P256-Y0708-01.webp
% 本文件只在转写完成后作为题干使用；OCR 条目尚不是可靠题干。
\begin{bankitem}{习题6.4 第10题}{第6章 / 6.4 用样本估计总体 / 教材 p.256 / PDF 第263页}{XJ-B1-C06-HW-P256-Y0708}{OCR未校对：公式与文字均可能有误；未完成数学TeX转写；禁止直接发给学生}
\begin{ocrdraft}
[教材p.256 / PDF第263页]
\par
10．一家超市想调查顾客的购买力，随机查阅了10位使用现金的顾客的消费额
\par
和12位使用移动支付的顾客的消费额，获得如下数据（单位：元）：
\par
现金消费：19．50
\par
39．50
\par
7\&60
\par
35．70
\par
31.80
\par
9\&30
\par
23．50
\par
10\&00
\par
29．00
\par
32．20
\par
移动支付消费：45．90
\par
89．70
\par
123．60
\par
98．75
\par
59．30
\par
45．35
\par
109．45
\par
95．50
\par
103．00
\par
2上00
\par
65．50
\par
67．20
\par
（1）分别估计该超市使用现金与使用移动支付的顾客的平均消费额；
\par
（2）在所有购物的顾客中，如果有32％的顾客使用现金，估计该超市的顾客的
\par
平均消费额．
\end{ocrdraft}
\end{bankitem}
